```
\begin{figure}[ht]
\centering
\begin{tikzpicture}[
  >={Stealth[scale=1.2]},
  node distance=1.5cm,
  box/.style={rectangle, draw, thick, minimum width=2cm, minimum height=1cm},
  circle/.style={circle, draw, thick, minimum size=1cm}
]
  % Nodes
  \node[box] (anchor) {Anchor};
  \node[box, right=of anchor] (positive) {Positive};
  \node[box, below=of positive] (negative1) {Negative 1};
  \node[box, below=of negative1] (negative2) {Negative 2};
  \node[circle, right=of positive] (loss) {Contrastive Loss};

  % Arrows
  \draw[->] (anchor) -- (positive) node[midway, above] {Similar};
  \draw[->] (anchor) -- (negative1) node[midway, above] {Dissimilar};
  \draw[->] (anchor) -- (negative2) node[midway, above] {Dissimilar};
  \draw[->] (positive) -- (loss);
  \draw[->] (negative1) -- (loss);
  \draw[->] (negative2) -- (loss);

  % Labels
  \node[above=0.1cm of anchor] {Input};
  \node[below=0.1cm of loss] {$\mathcal{L} = \log \frac{e^{\text{sim}(a,p)}}{e^{\text{sim}(a,p)} + \sum_k e^{\text{sim}(a,n_k)}}$};
\end{tikzpicture}
\caption{Contrastive loss structure with anchor, positive, and negative samples.}
\label{fig:fig_contrastive_loss}
\end{figure}
```